% Candidate A.55. Notes only; not registered before independent audit.
\subsection{Cyclic auxiliary covers at the coefficient prime}

\begin{statement}{A.55}{Proved}{All-prime normalizers for the quartic-CM genus-two target and the Heisenberg-five cover}
Let
\[
 G:\ v^2=w^5-1,
 \qquad G^\circ=G\setminus\{(0,i),(0,-i),\infty\}.
                                                        \tag{A.55.1}
\]
Over a number field defining this curve, for every prime \(\ell\),
every nontrivial closed procyclic subgroup of
\(\pi_1(G^\circ_{\overline K})^{(\ell)}\) whose arithmetic
normalizer has open image is contained in a conjugate of cusp inertia.
The assertion is unchanged by a finite field extension or a twist
which becomes isomorphic to this punctured curve after such an extension.

Consequently the normalizer assertion of A.54 holds at \emph{every}
coefficient prime, including all \(\ell\equiv1\pmod5\). In particular
it holds at \(\ell=11\) for the actual genus-51, degree-125
Heisenberg cover in A.54. This statement adds no fixed-subgroup
claim to the pro-five fixed-subgroup result already proved there.
\end{statement}
\begin{proof}
\paragraph{The integral cyclic-cover calculation.}
We first record precisely the integral detector that will be used.
Let \(q:Y\to X\) be the connected unramified cyclic cover of
degree \(m=\ell^n\) of a proper genus-two curve defined by a
primitive character
\(\eta:H_1(X,\Z_\ell)\to\Z/m\).
Choose a symplectic basis \(a,b,c,d\) so that
\(\eta(b)=1\) and \(\eta(a)=\eta(c)=\eta(d)=0\).
The line generated by \(a\) modulo \(m\) is the Poincare-dual
line of \(\eta\), with an immaterial sign convention.
Then
\[
 q_*\bigl(H_1(Y,\Z_\ell)^{\operatorname{Deck}(q)}\bigr)
   =\Z_\ell a\oplus m\Z_\ell b
       \oplus m\Z_\ell c\oplus m\Z_\ell d.             \tag{A.55.2}
\]
Here and below these homology groups can equivalently be taken as
the integral Tate homology of the respective Jacobians.

For completeness this is an integral surface-group calculation.
Use the presentation
\(\langle a,b,c,d\mid[a,b][c,d]=1\rangle\) and
coset representatives \(1,b,\ldots,b^{m-1}\).
The Schreier generators are
\[
 a_j=b^jab^{-j},\quad c_j=b^jcb^{-j},\quad
 d_j=b^jdb^{-j}\quad(0\le j<m),\qquad B=b^m.
\]
After abelianization the \(m\) lifted surface relations are
\(a_j-a_{j+1}=0\), with indices taken cyclically. Their relation
lattice is saturated: eliminate \(a_1,\ldots,a_{m-1}\), leaving
the single generator \(h=[a_0]\). Thus an integral basis of
\(H_1(Y,\Z_\ell)\) is
\(h,B,[c_j],[d_j]\), of rank \(2m+2\).
The deck generator fixes \(h,B\) and cyclically permutes each
of the two \(m\)-tuples. Its invariant lattice has basis
\(h,B,\sum_j[c_j],\sum_j[d_j]\). Their pushforwards are
\(a,mb,mc,md\), proving (A.55.2).
In particular, a class whose rational image is in the old subspace
\(q^*H_1(X,\Q_\ell)\) satisfies
\[
      q_*[u]\bmod m\ \in\ (\Z/m)\operatorname{PD}(\eta).
                                                        \tag{A.55.3}
\]
Indeed the old rational subspace is exactly the deck-invariant
subspace, and a deck-invariant rational integral vector is
integrally invariant because the integral homology is torsion-free.
This implication is integral; it does not assume that \([u]\) or
\(q_*[u]\) is primitive.

\paragraph{The local auxiliary covers and their Pryms.}
We next prove the local geometric input, for arbitrary \(n\).
Let \(\mathcal X/R\) be a smooth proper curve over a complete
mixed-characteristic discrete valuation ring of residue
characteristic \(p\), and suppose a degree-zero line bundle
\(\mathcal L\) has order \(m=p^n\), with special line bundle
\(\overline{\mathcal L}\) still of exact order \(m\).
After a finite extension choose an isomorphism
\(\mathcal L^{\otimes m}\simeq\mathcal O\). It defines a finite
flat \(\mu_m\)-torsor
\[
 \mathcal Y=\operatorname{Spec}_{\mathcal X}
        \left(\bigoplus_{j=0}^{m-1}\mathcal L^{\otimes j}\right)
                         \longrightarrow\mathcal X,       \tag{A.55.4}
\]
where the multiplication uses this trivialization; reversing
\(\mathcal L\) changes only the convention. Its generic fibre
is a connected etale cyclic cover of degree \(m\).

The special fibre of (A.55.4) is integral and its normalized
function field is the purely inseparable extension
\(k(X)^{1/p^n}\) of \(k(X)\). To see both assertions, represent
\(\overline{\mathcal L}\) by a divisor \(D\) and its trivialization
by \(f\) with \(\operatorname{div}(f)=mD\). If \(f=g^p\),
then \(p^{n-1}D=\operatorname{div}(g)\), contrary to the exact
order of \(\overline{\mathcal L}\). Hence \(f\notin k(X)^p\),
and adjoining an \(m\)-th root gives degree \(m\).
A one-variable function field over a perfect field has
\([k(X)^{1/p^n}:k(X)]=p^n\), so the resulting field is exactly
the displayed Frobenius extension. Finite flatness makes each
special affine ring torsion-free over the base affine ring;
its injection into this generic field proves that it is a domain.

The total space \(\mathcal Y\) is normal. Locally it is the
hypersurface \(z^m=u\), with \(u\) a unit, over a regular
two-dimensional base, so it satisfies \(S_2\). Horizontal
codimension-one points lie in the smooth generic fibre. At the
unique vertical generic point the special fibre is generically
reduced, so the local maximal ideal is generated by a base
uniformizer and the local ring is regular. This proves \(R_1\),
hence normality. The same reasoning applies after every finite
extension of the valuation ring, since the special order of
\(\overline{\mathcal L}\) is unchanged.

Apply \citet[Theorem 1, printed p.179]{fghRaynaud1990Pgroups}
to this actual unramified \(p\)-group cover of the generic fibre
of the smooth proper model. Its Jacobian has potentially good
reduction. In a semistable model the graph is a tree, and the
components which map to points of the original special fibre
have crystalline slopes strictly between zero and one.
There is exactly one nonexceptional component: it is the
normalization \(X^{(p^{-n})}\) just identified, with its radicial
degree-\(m\) map to \(X\). The reduced special fibre is
geometrically unibranch because this normalization is radicial;
thus there is no additional nonexceptional component hidden at
a singular point.

Let \(P_q=(\ker(q_*:J_Y\to J_X))^0\) be the Prym.
After taking good models, the norm extends as a homomorphism of
abelian schemes. The special Jacobian of the compact-type fibre
is the product of its component Jacobians. On the unique main
factor the special norm is the isogeny induced by the radicial
Frobenius map; on every exceptional factor it is zero, because
that component maps to a point. It follows, up to isogeny, that
\[
        \overline{P_q}\ \sim\
           \prod_{Z\ \mathrm{exceptional}}\operatorname{Jac}(Z).
                                                        \tag{A.55.5}
\]
This can also be read from the rational idempotent
\(1-m^{-1}q^*q_*\), which extends after multiplication by \(m\).
In particular, \(P_q\) has potentially good reduction and
special \(p\)-rank zero. The main factor need not have
\(p\)-rank zero; it is precisely the old factor removed in
(A.55.5).

The rank-one crystalline argument of A.46 now gives
\[
 V_p(P_q)\text{ has no nonzero }\Q_p\text{-line stable under
 any open local or global Galois subgroup}.              \tag{A.55.6}
\]
Indeed, dualizing a stable homological line gives a rank-one
crystalline quotient of \(H^1\). In the cohomological convention
its Hodge--Tate weights are zero or one; the rank-one
weak-admissibility argument of A.46 would give Newton slope
zero or one. Neither slope occurs in (A.55.5). This conclusion is unchanged
by further finite extensions. Notice that Raynaud's theorem is
used for the full cyclic group of order \(p^n\); a calculation
only for degree \(p\) would not suffice below.

\paragraph{The split coefficient primes on the actual genus-two curve.}
The Jacobian of \(G\) has full CM by \(\Z[\zeta_5]\), with
type \(\{1,2\}\). A.54 proves that its four labelled
one-dimensional CM characters remain distinct on every open
Galois subgroup. At a nonsplit coefficient prime there is no
\(\Q_\ell\)-line stable under an open subgroup, by that same
argument, so A.11 already proves the assertion.

It remains to consider \(p=\ell\equiv1\pmod5\).
After a finite field extension defining all endomorphisms, the
integral Tate lattice has four rank-one eigensummands for
\(\Z_p[\zeta_5]\simeq\Z_p^4\). They form the two perfect
symplectic pairs with labels \(\{1,4\}\) and \(\{2,3\}\).
The integral decomposition is valid because \(p\ne5\).
The curve has good ordinary reduction at a place over \(p\).
For a direct check, write \(p=10k+1\). On the basis
\(dw/v,w\,dw/v\), the Cartier matrix of \(v^2=w^5-1\)
over \(\F_p\) is diagonal, with nonzero diagonal entries
\[
        (-1)^{3k}\binom{5k}{2k},\qquad
        (-1)^k\binom{5k}{4k}.                         \tag{A.55.7}
\]
Indeed its entries are the coefficients of \(w^{pi-j}\)
in \((w^5-1)^{(p-1)/2}\), with indices \(1\le i,j\le2\);
the off-diagonal exponents are not divisible by five. The
binomial coefficients are units since \(5k<p\).
The coefficient formula follows directly by applying Cartier
to \(w^{j-1}v^{-p}(w^5-1)^{(p-1)/2}dw\), as in
\cite[Section 3, pp.5--6]{fghAchterHowe2017}.
Invertibility is the ordinary criterion.

The four CM idempotents extend to the good abelian scheme and
split its \(p\)-divisible group into height-one factors.
Polarization makes the two factors in each symplectic pair
Cartier dual; consequently exactly one in each pair is etale
and the other multiplicative. In particular, in either chosen
pair there are compatible Picard torsion points \(L_n\) of
order \(p^n\) whose specializations have exactly that order.
They are defined over finite number fields at each finite
level. Their line bundles extend over the good curve model:
use the relative degree-zero Picard scheme and rigidify at a
rational point, after another finite extension. The trivialization
\(L_n^{\otimes p^n}\simeq\mathcal O\) also extends after
rigidification; any scalar discrepancy is absorbed after a
finite extension. Thus these are the actual line bundles
required in (A.55.4), not merely classes in a local vector space.

Suppose a noncuspidal procyclic subgroup \(A\) of the
geometric pro-\(p\) group of \(G^\circ\) has open arithmetic
normalizer. Choose a generator \(u\). By A.11 its proper
homology class \(v\) is nonzero, and its rational span is
stable under an open Galois subgroup. Distinctness of the CM
characters puts \(v\) in a single one of the four integral
eigenlines. Write \(v=p^r v_0\), where \(v_0\) is primitive.
Choose the etale Picard eigenline in the \emph{other} symplectic
pair and take \(n>r\).

Kummer duality and the principal polarization identify
\(L_n\in J_G[p^n]\) with the character
\(\eta_n:H_1(G,\Z_p)\to\mu_{p^n}\) obtained by Weil pairing.
After fixing a primitive root, identify its target with
\(\Z/p^n\). Its Poincare-dual line is the chosen Picard
line, up to a unit and the harmless choice of signs. Hence
\[
       \eta_n(v)=0,
       \qquad
       v\bmod p^n\notin
          (\Z/p^n)\operatorname{PD}(\eta_n).          \tag{A.55.8}
\]
The second assertion uses \(n>r\) and the direct integral
decomposition into the two different symplectic pairs. This
is why the other pair was chosen; choosing the same dual line
could give no contradiction to (A.55.3).

Let \(q_n:Y_n\to G\) be the actual cyclic unramified proper
cover supplied by this Kummer torsor, and set
\(Y_n^\circ=q_n^{-1}(G^\circ)\). The kernel of \(\eta_n\)
contains \(A\); the geometric pro-\(p\) group of
\(Y_n^\circ\) is that open index-\(p^n\) subgroup, because
the pro-\(p\) topology on an open subgroup is induced.
Choose the corresponding lift of the base point. Descend the
cover and this point to a finite field extension and intersect
the original normalizer with its arithmetic fundamental group.
This intersection still has open arithmetic image and still
normalizes \(A\). Thus the proper class
\(\widetilde v\in H_1(Y_n,\Z_p)\) of the \emph{same}
generator \(u\) spans a stable line, with
\(q_{n,*}\widetilde v=v\ne0\). These are actual integral
classes in the covering group, not rational Malcev lifts.

The projection of this line to \(V_p(P_{q_n})\) must vanish
by (A.55.6). Therefore \(\widetilde v\) is in the old rational
homology, and (A.55.3) gives the opposite of (A.55.8).
This contradiction proves the normalizer criterion at every
split prime. The choice of \(n\) also handles arbitrary
divisibility of the original proper class.

\paragraph{Return to the Heisenberg cover.}
For the actual cover of A.54, its proper homology is
\(V_\ell E_i\oplus V_\ell A_2^{\oplus5}
\oplus V_\ell B_8^{\oplus5}\).
The \(B_8\) summand has no potentially stable rational line,
as proved there. A line projecting nontrivially to \(E_i\)
is excluded by the original punctured elliptic covering map.
The remaining lines are detected by the deck translates of
the actual maps to \(D_\lambda/\sigma\) constructed in A.54;
their proper targets are twists of \(G\), and their removed
fibres are exactly the three points in (A.55.1).
Indeed, after a finite scaling, the target is
\(u^5=t(t-1)\), related to (A.55.1) by
\(t=(v+i)/(2i)\) and \(w=\beta u\), \(\beta^5=-4\).
The maps are therefore maps of the actual punctured curves.
Their nonzero images would have open arithmetic normalizer
in the target pro-\(\ell\) group and nonzero proper homology,
contradicting the criterion just proved. Statement A.11
finishes the all-prime assertion for the Heisenberg cover.

The exact certificate accompanying this proof checks the
Schreier relations, integral deck invariant lattice and its
pushforward for prime-power degrees, and (A.55.7) at all
split primes below a fixed bound. The uniform statements are
proved above and are not inferred from the finite checks.
\end{proof}
